\section{Counting complete runs}\label{app:complete-runs}

The complete-run bound counts a whole run at once, coupling free returns
with returns that add a zero equation. It is uniform over the active
forms and can be used with the Newton or subset-sum precode.

\paragraph{A bound for all free returns.}
Let $\mu$ be a distribution on nonzero vectors of a finite-dimensional
space over $\F_p$, with $p$ odd. Draw independent increments with law
$\mu$, set $S_0=0$, write $S_j$ for their partial sum, and let $T$ be the first
positive return time to zero. For $0<\omega<1$, define
\[
 G_\mu(\omega)=\sum_{j\ge0}\Pr[S_j=0]\omega^j,
 \qquad F_\mu(\omega)=\E[\omega^T;T<\infty].
\]
The renewal identity gives $F_\mu=1-1/G_\mu$.
Uniformly over $\mu$, including nonsymmetric distributions,
\begin{equation}\label{eq:complete-free-return}
 F_\mu(\omega)\le F_p(\omega):=1-\frac1{G_p(\omega)},\qquad
 G_p(\omega)=\frac1p\sum_{j=0}^{p-1}\frac1{1-\omega\cos(2\pi j/p)}.
\end{equation}
To prove this, average over additive characters $\chi$ of the vector
space and put $\phi_\chi=\sum_a\mu(a)\chi(a)$. Fourier inversion gives
\[
 \begin{aligned}
 G_\mu(\omega)
 &=\E_\chi\operatorname{Re}\frac1{1-\omega\phi_\chi}
 \le\E_\chi\frac1{1-\omega\operatorname{Re}\phi_\chi}\\
 &\le\sum_a\mu(a)\E_\chi\frac1{1-\omega\operatorname{Re}\chi(a)}
 =G_p(\omega).
 \end{aligned}
\]
The second inequality is Jensen's inequality for the convex function
$u\mapsto(1-\omega u)^{-1}$. Every nonzero $a$ has order $p$, so its character
values are uniform over the $p$th roots of unity. Balanced increments
$a,-a$ attain the bound. The Poisson-kernel series also gives
\begin{equation}\label{eq:finite-cycle-green}
 G_p(\omega)=\frac1{\sqrt{1-\omega^2}}\frac{1+r_\omega^p}{1-r_\omega^p},
 \qquad r_\omega=\frac{\omega}{1+\sqrt{1-\omega^2}}.
\end{equation}
Thus $0\le F_p(\omega)-(1-\sqrt{1-\omega^2})\le2r_\omega^p$. The term
$1-\sqrt{1-\omega^2}$ is the $p$-independent main term of $F_p(\omega)$.

\paragraph{Returns and rank gains in the same run.}
Fix a run start, its exposure history, and its quotient by earlier zero
equations. Let $a_1,\ldots,a_m$ be the forms of the remaining active
input labels in this quotient. On a retained branch every $a_i$ is
nonzero. Define
\[
 \mu=\frac1m\sum_{i=1}^m\delta_{a_i},\qquad
 \mathcal U=\bigcup_{i=1}^m
                 (\operatorname{span}(a_i)\setminus\{0\}).
\]
Here $\delta_a$ denotes point mass at $a$. Until the next zero output,
the quotient remains fixed. All these input forms remain required
nonzero, including labels consumed during the run. Imposing a zero
equation on a nonzero sum in $\mathcal U$ would make some required form
$a_i$ zero, so no such return occurs on a retained branch.
Outside $\mathcal U$, allow a rank-gaining return with
weight $t=1/\tau$, where $\tau\ge1$. Ignoring further nonzero requirements
can only enlarge the count. For independent draws from the fixed
distribution $\mu$, the complete-run series is
\[
 \mathcal R_\mu(\omega,t)=F_\mu(\omega)
       +t\sum_{j\ge1}\Pr[T>j,\ S_j\notin\mathcal U]\omega^j.
\]
The event $T>j$ excludes the zero vector in the second term and every
earlier free return. This is a bound on a weighted branch count, not
the probability of one event: a walk may allow rank-gaining endpoints
at several lengths. Distinct-label sampling is restored below.

Assume $0<\omega<1$ and $\omega\le c\le1$, where $c=t/(1-\omega)$.
For a one-dimensional subspace $\ell$ with $\mu(\ell)>0$,
write $a_\ell=\mu(\ell)$, $\theta_\ell=\omega a_\ell$, and let $\nu_\ell$ be
the conditional increment distribution on that line. Starting from
$f\in\ell\setminus\{0\}$, let $H_\ell(f)$ be the $\theta_\ell$-discounted
probability of hitting zero under $\nu_\ell$. A hit needs at least one
step, so $0\le H_\ell(f)\le\theta_\ell$. Put
\[
 b_\ell=\frac{c\omega(1-a_\ell)}{1-\omega a_\ell},\qquad
 V(f)=\begin{cases}
 b_\ell+(1-b_\ell)H_\ell(f),&f\in\ell\setminus\{0\}\subset\mathcal U,\\
 c,&f\ne0,\ f\notin\mathcal U.
 \end{cases}
\]
On an active line,
$V(f)\le\omega[c+(1-c)a_\ell]\le\omega\le c$.
Outside $\mathcal U$, no active increment cancels $f$; its immediate
rank-gaining return and subsequent continuations are bounded by
$t+\omega c=c$. On a line, off-line increments have continuation at most $c$,
and on-line transitions satisfy the hitting-probability recurrence.
Substitution gives exactly $V(f)$ as their upper bound. Induction over
finite truncations, followed by monotone convergence, proves
\[
 \mathcal R_\mu(\omega,t)
 \le\sum_\ell\left[\theta_\ell b_\ell+(1-b_\ell)F_p(\theta_\ell)\right].
\]
Here $F_{\nu_\ell}(\theta_\ell)=\theta_\ell\E_{a\gets\nu_\ell}H_\ell(a)$,
and~\eqref{eq:complete-free-return} bounds each line term.

Since $\sum_\ell\theta_\ell=\omega$, bound each $F_p(\theta_\ell)$ by its
main term $1-\sqrt{1-\theta_\ell^2}$ plus $2r_{\theta_\ell}^p$.
Substituting $\theta_\ell=2r/(1+r^2)$ into the main terms gives
\begin{equation}\label{eq:complete-line-envelope}
 \mathcal R_\mu(\omega,t)\le\omega\max_{0\le r\le r_\omega}\varphi(r)
       +2r_\omega^p,
 \quad
 \varphi(r)=(1-c\omega)r+c(2-\omega)-\frac{2t}{1-r}.
\end{equation}
The finite-field error follows from $r_{\theta_\ell}/r_\omega\le\theta_\ell/\omega$
and $\sum_\ell(\theta_\ell/\omega)^p\le1$. This bound retains the dependence
between free returns and rank gains; it does not maximize their
probabilities separately.

\paragraph{Distinct labels and terminal words.}
Choose $x>0$, $0<y\le1$, $\lambda>1$, and set
\[
 A=\frac{xy}{\lambda-1},\qquad B=\frac{x}{\lambda-y},\qquad \omega=By<1.
\]
After summing zero-input gaps, a complete run with $j$ active
inputs has factor $AB^{j-1}y^{j-2}=(A/(By^2))\,\omega^j$. The
run-series multiplier $A/(By^2)$ therefore converts the complete-run
series, whose $j$-input terms carry $\omega^j$, into these run factors.
Use cutoff $J=16$ and budget $\beta(\kappa)=\max(1,\kappa^J)$. Write
$(\kappa)_j=\kappa(\kappa-1)\cdots(\kappa-j+1)$ for the falling factorial.
For $\kappa>J$ and $j\le J$, conditioning independent labels to be distinct
is absorbed by
\[
 \frac{\kappa^j}{(\kappa)_j}\beta(\kappa-j)\le\beta(\kappa).
\]
Indeed, after dividing by $\kappa^J$ the left side is
$((\kappa-j)^j/(\kappa)_j)((\kappa-j)/\kappa)^{J-j}\le1$.
For $j>J$, work directly without replacement. Free and rank-gaining
endpoints are disjoint and have total rank-weighted probability at most
one. Their geometric tail is at most $\omega^{J+1}/(1-\omega)$.

The finite certificate requires $\omega\le1/2$ and $p\ge256$, so the
finite-field term in~\eqref{eq:complete-line-envelope} is at most $2^{-255}$.
With slack $\epsilon=10^{-6}$, it suffices to check
\begin{equation}\label{eq:complete-run-witness}
 \frac{A}{By^2}\left[\omega\max_{0\le r\le r_\omega}\varphi(r)
             +2^{-255}+\frac{\omega^{17}}{1-\omega}\right]\le1-\epsilon.
\end{equation}
The extra tail deliberately overcounts long runs already present in the
independent-draw envelope. Terminal words and pools of at most sixteen
active labels are bounded by
\begin{align*}
 R_{\rm terminal}&=\frac\lambda{\lambda-1}
                 +\frac{A\lambda}{(\lambda-y)(1-\omega)},\\
 R_{\rm small}&=\frac\lambda{\lambda-1}
       \sum_{j=0}^{16}\left(\frac{x(1+y)}{\lambda-1}\right)^j,
 \qquad C=\frac{\max(R_{\rm terminal},R_{\rm small})}\epsilon.
\end{align*}
To make the induction explicit, discount a word of length $N$
by $\lambda^{-N}$. Its marker weight is $x^\kappa y^h$, where $\kappa,h$ are
the input and output weights, and each new zero equation contributes
$\tau^{-1}$. Average over the uniform order of the active labels and
sum over support branches and placements, as in
Lemma~\ref{lem:scan-composition}.

For $\kappa\le16$, $R_{\rm small}\le\epsilon C$ bounds the entire residual
word, after summing its possible lengths. For $\kappa>16$, split
at the first completed run. The distinct-label inequalities
and~\eqref{eq:complete-run-witness} bound the completed-run contribution
by $(1-\epsilon)C\beta(\kappa)$, assuming the bound for smaller pools.
The terminal contribution is at most $R_{\rm terminal}\le\epsilon C$.
Since $\beta(\kappa)\ge1$, their sum is at most $C\beta(\kappa)$.
The bound is uniform over the new quotient after the return, which
justifies its use in the induction.

Extracting length $n$ gives
$K(\kappa)=C\beta(\kappa)\lambda^n$, including either final feedback state.
For positive $s,h$, Lemma~\ref{lem:scan-composition} therefore gives
the factor $(sh)^{16}$ multiplying~\eqref{eq:width-one-profile}.
The source factor remains $F_s(\tau)/\tau$; it is charged only once.

\paragraph{Rational checks and the Newton instantiation.}
The certificate specifies, for each use of~\eqref{eq:complete-run-witness},
rational $\tau,x,y,\lambda$ and $0<\rho<1$ satisfying
$\omega(1+\rho^2)\le2\rho$. This implies $\rho\ge r_\omega$.
After subtracting the explicit error and tail terms,
inequality~\eqref{eq:complete-run-witness} reduces to $\varphi(r)\le H$ on
$[0,\rho]$, for a computed rational $H$. Multiplying by $1-r$ gives
\[
 (c\omega-H)+(1-2c+H)r+(c\omega-1)r^2\le0.
\]
Its leading coefficient is negative. Exact evaluation at both endpoints
and at its vertex, when inside the interval, proves the required
uniform bound. All checks are exact rational comparisons.

For the Newton precode, retain $Q_8(z,\tau)$
from~\eqref{eq:newton-source-polynomial}. At forbidden weight
$e=245366$ (so the distance is at least $e+1$), the hypergeometric
bound~\eqref{eq:hyper-sparse} covers all $s<32$. This sparse
contribution has base-two logarithm at most $-49.647382836382$.
The dense domain is partitioned recursively into rectangles; each cell
is a leaf. There are 180 nonempty leaves: 44 use the projective-pair
bound, three use the joint-renewal bound, and 133 use the complete-run
bound. Every integer profile with $32\le s\le\min(n,4e)$ and
$\lceil s/2\rceil\le h\le\min(n,2e)$ is covered. The total failure
probability has base-two logarithm at most $-49.647365460353$. Thus the
unchanged addition-only Newton encoder has distance at least $245367$
($>11.7\%$ of $n$) with failure probability below $2^{-49.64}$.
The certificate and its verifier are listed in
Appendix~\ref{app:reproduction}.

The argument is uniform in the source forms. Substituting the catalog
of Appendix~\ref{app:subset-sum-gv} gives the stronger distance of
Theorem~\ref{thm:subset-sum-gv}, with a smaller failure margin and
fewer additions.
